% Oefeningen bij dit hoofdstuk (ingevoegd achteraan het hoofdstuk met \input).
% Een oefening:  \begin{oefening}[titel] ... \end{oefening}   (boek en website)
% De oplossing:  \begin{oplossing} ... \end{oplossing}       (enkel website, dichtgeklapt)

Neem $N_A = \SI{6.022e23}{\per\mole}$, $k_B = \SI{8.617e-5}{\electronvolt\per\kelvin}$, $R = \SI{8.314}{\joule\per\mole\per\kelvin}$, $F = \SI{96485}{\coulomb\per\mole}$, $hc = \SI{1240}{\electronvolt\nano\metre}$ en voor Cu K$\alpha$-straling $\lambda = \SI{0.1541}{\nano\metre}$.

\begin{oefening}[Waarom is aluminiumfolie zo licht? Dichtheid uit de kristalstructuur]
IJzer (\(\alpha\)-Fe, \(M = 55.85\,\mathrm{g/mol}\)) kristalliseert bij kamertemperatuur in een BCC-rooster met \(a = 0.2866\,\mathrm{nm}\). Aluminium (\(M = 26.98\,\mathrm{g/mol}\)) kristalliseert in een FCC-rooster met \(a = 0.405\,\mathrm{nm}\).

\begin{enumerate}
\item[(a)]
  In een BCC-rooster raken de atomen langs de lichaamsdiagonaal. Bereken de atoomstraal van ijzer en de pakkingsfractie van het BCC-rooster.
\item[(b)]
  Bereken de dichtheid van ijzer en van aluminium uit de kristalgegevens. Vergelijk met de tabelwaarden (\(7.87\) en \(2.70\,\mathrm{g/cm^3}\)).
\item[(c)]
  Aluminium heeft de dichtste pakking, maar is toch bijna drie keer lichter dan ijzer. Waarom?
\item[(d)]
  Bereken de pakkingsfractie van een primitief kubisch rooster. Waarom komt dat rooster bij metalen bijna niet voor?
\end{enumerate}
\end{oefening}
\begin{oplossing}
\textbf{Resultaat:} (a) \(r_{\mathrm{Fe}} =\) \textbf{0.124 nm}, pakkingsfractie \textbf{0.68}. (b) \(\rho_{\mathrm{Fe}} =\) \textbf{7.88 g/cm\ensuremath{^{3}}}, \(\rho_{\mathrm{Al}} =\) \textbf{2.70 g/cm\ensuremath{^{3}}}, in uitstekende overeenstemming. (c) Een Al-atoom is \textbf{half zo zwaar} en \textbf{groter}. (d) \(\pi/6 =\) \textbf{0.52}: te veel lege ruimte.

\textbf{(a)} De lichaamsdiagonaal heeft lengte \(a\sqrt{3}\) en bevat \(4r\) (de straal van elk van de twee hoekatomen plus de diameter van het centrale atoom): \(r = a\sqrt{3}/4 = 0.2866 \times 1.732/4 = 0.124\,\mathrm{nm}\). De cel bevat 2 atomen:

\[
f = \frac{2\cdot\frac{4}{3}\pi r^3}{a^3} = \frac{\pi\sqrt{3}}{8} = 0.68.
\]

\textbf{(b)} \(\rho = \dfrac{Z\,M}{N_A a^3}\) met \(Z\) het aantal atomen per cel:

\begin{itemize}
\item
  Fe: \(\rho = \dfrac{2 \times 55.85}{(6.022\times10^{23})(0.2866\times10^{-7}\,\mathrm{cm})^3} = 7.88\,\mathrm{g/cm^3}\)
\item
  Al: \(\rho = \dfrac{4 \times 26.98}{(6.022\times10^{23})(0.405\times10^{-7}\,\mathrm{cm})^3} = 2.70\,\mathrm{g/cm^3}\)
\end{itemize}

Dit is een mooie controle van een XRD-meting: de dichtheid van een kristal volgt rechtstreeks uit de roosterconstante.

\textbf{(c)} De pakking van Al (0.74) is maar 9 \% hoger dan die van Fe (0.68). Maar een Al-atoom weegt minder dan de helft van een Fe-atoom en is bovendien groter (\(r = a\sqrt{2}/4 = 0.143\,\mathrm{nm}\) tegenover \(0.124\,\mathrm{nm}\)). De dichtheid hangt dus vooral af van de atoommassa en -grootte, en veel minder van het roostertype. Daarom bouwt men vliegtuigen en drankblikjes uit aluminium.

\textbf{(d)} Eén atoom per cel met \(2r = a\): \(f = \frac{4}{3}\pi(a/2)^3/a^3 = \pi/6 = 0.52\). Bijna de helft is lege ruimte; bij niet-gerichte metaalbindingen levert een dichtere pakking meer buren en dus een lagere energie. Enkel polonium kristalliseert in dit rooster.
\end{oplossing}

\begin{oefening}[IJzer krimpt bij verhitten: de \(\alpha \rightarrow \gamma\)-overgang]
Bij \(912\,^\circ\mathrm{C}\) gaat BCC-ijzer (\(a = 0.2904\,\mathrm{nm}\)) over in FCC-ijzer (\(a = 0.3647\,\mathrm{nm}\), beide roosterconstanten bij de overgangstemperatuur). De overgangsenthalpie bedraagt ongeveer \(+0.90\,\mathrm{kJ/mol}\).

\begin{enumerate}
\item[(a)]
  Bereken het volume per atoom in beide structuren. Zet het metaal uit of krimpt het bij de overgang?
\item[(b)]
  Bereken \(\Delta V_m\) en gebruik de vergelijking van Clapeyron (hoofdstuk 5), \(\mathrm{d}T/\mathrm{d}p = T\Delta V_m/\Delta H_m\), om te voorspellen hoe de overgangstemperatuur verandert met de druk.
\item[(c)]
  Een smid verhit een stuk staal tot het gloeit en dompelt het dan in water. Waarom kunnen er bij zulke snelle overgangen tussen kristalstructuren scheurtjes ontstaan?
\end{enumerate}
\end{oefening}
\begin{oplossing}
\textbf{Resultaat:} (a) \(0.01225\,\mathrm{nm^3}\) (BCC) tegenover \(0.01213\,\mathrm{nm^3}\) (FCC): ijzer \textbf{krimpt ongeveer 1 \%}. (b) \(\Delta V_m = -7.1\times10^{-8}\,\mathrm{m^3/mol}\), \(\mathrm{d}T/\mathrm{d}p \approx\) \textbf{\ensuremath{-}9 K per kbar}. (c) \textbf{Volumesprongen} in een niet-uniform afgekoeld stuk veroorzaken inwendige spanningen.

\textbf{(a)} BCC: 2 atomen per cel, \(V = a^3/2 = (0.2904)^3/2 = 0.01225\,\mathrm{nm^3}\). FCC: 4 atomen per cel, \(V = a^3/4 = (0.3647)^3/4 = 0.01213\,\mathrm{nm^3}\). Het volume daalt met \(0.97\,\%\): ijzer krimpt bij verhitten door de overgang, omdat FCC dichter gepakt is (0.74 tegenover 0.68). Met een dilatometer zie je die krimp als een duidelijke knik in de uitzettingscurve.

\textbf{(b)} \(\Delta V_m = (0.01213 - 0.01225)\times10^{-27}\,\mathrm{m^3} \times N_A = -7.1\times10^{-8}\,\mathrm{m^3/mol}\). Met \(T = 1185\,\mathrm{K}\):

\[
\frac{\mathrm{d}T}{\mathrm{d}p} = \frac{(1185)(-7.1\times10^{-8})}{900} = -9.4\times10^{-8}\,\mathrm{K/Pa} \approx -9\,\mathrm{K/kbar}.
\]

Druk bevoordeelt de dichtere FCC-fase, zodat de overgang bij hogere druk bij een lagere temperatuur plaatsvindt (Le Châtelier). Het effect is klein, omdat \(\Delta V\) in de vaste fase klein is (§10.6.1). In de aardkern, bij miljoenen bar, speelt het wel een grote rol.

\textbf{(c)} Bij snel afkoelen koelt de buitenkant eerst en verandert daar eerst van structuur, met een volumesprong, terwijl de kern nog in de andere structuur zit. Dat geeft grote inwendige spanningen. In koolstofstaal ontstaat bij afschrikken bovendien martensiet, een vervormde structuur waarin koolstof gevangen zit en die de dislocaties blokkeert: het staal wordt keihard maar bros (§10.7). Daarom wordt het geharde staal achteraf ``ontlaten'' bij een gematigde temperatuur.
\end{oplossing}

\begin{oefening}[Een korrel keukenzout]
Keukenzout heeft de NaCl-structuur (§10.2.3) met \(a = 0.564\,\mathrm{nm}\) (\(M = 58.44\,\mathrm{g/mol}\)).

\begin{enumerate}
\item[(a)]
  Hoeveel \(\mathrm{Na^+}\)- en \(\mathrm{Cl^-}\)-ionen bevat een elementaire cel? Toon dit aan door te tellen.
\item[(b)]
  Bereken de dichtheid van NaCl.
\item[(c)]
  Wat is de kortste afstand tussen een \(\mathrm{Na^+}\)- en een \(\mathrm{Cl^-}\)-ion? Vergelijk met de som van de ionstralen (\(0.102 + 0.181\,\mathrm{nm}\)).
\item[(d)]
  Een zoutkorrel is een kubusje met een ribbe van \(0.5\,\mathrm{mm}\). Hoeveel elementaire cellen en hoeveel formule-eenheden bevat het?
\item[(e)]
  Waarom hebben zoutkorrels onder een loep de vorm van kubusjes?
\end{enumerate}
\end{oefening}
\begin{oplossing}
\textbf{Resultaat:} (a) \textbf{4} \(\mathrm{Na^+}\) en \textbf{4} \(\mathrm{Cl^-}\). (b) \textbf{2.16 g/cm\ensuremath{^{3}}}. (c) \(a/2 =\) \textbf{0.282 nm}, gelijk aan de som van de ionstralen (\(0.283\,\mathrm{nm}\)). (d) \(7.0\times10^{17}\) cellen, \textbf{\(2.8\times10^{18}\)}** formule-eenheden**. (e) De kubische cel herhaalt zich langs drie loodrechte richtingen.

\textbf{(a)} \(\mathrm{Cl^-}\) vormt een FCC-rooster: \(8 \times \tfrac{1}{8} + 6 \times \tfrac{1}{2} = 4\). \(\mathrm{Na^+}\) zit in de oktaëdrische holtes: één in het centrum van de cel en één in het midden van elk van de 12 ribben (elk gedeeld door 4 cellen): \(1 + 12 \times \tfrac{1}{4} = 4\). Er zijn dus 4 formule-eenheden NaCl per cel.

\textbf{(b)} \(\rho = \dfrac{4 \times 58.44}{(6.022\times10^{23})(0.564\times10^{-7})^3} = 2.16\,\mathrm{g/cm^3}\).

\textbf{(c)} Een \(\mathrm{Na^+}\) in het midden van een ribbe ligt op \(a/2 = 0.282\,\mathrm{nm}\) van de \(\mathrm{Cl^-}\)-ionen op de hoeken. Dat komt overeen met \(0.102 + 0.181 = 0.283\,\mathrm{nm}\): de ionen raken elkaar langs de ribbe, zoals verwacht voor een ionaire binding (hoofdstuk 2).

\textbf{(d)} \((0.5\times10^{-3}/0.564\times10^{-9})^3 = (8.9\times10^{5})^3 = 7.0\times10^{17}\) cellen, dus \(2.8\times10^{18}\) formule-eenheden, of \(4.6\times10^{-6}\,\mathrm{mol}\) (\(0.27\,\mathrm{mg}\)). Een korrel bevat langs één ribbe bijna een miljoen cellen op een rij.

\textbf{(e)} Een kristal groeit door het aanhechten van ionen op de vlakken. De macroscopische vorm weerspiegelt de symmetrie van de elementaire cel: bij NaCl groeien de \((100)\)-vlakken het traagst en blijven ze dus als buitenvlakken over. Het kristal splijt ook netjes langs die vlakken, die enkel uit een afwisseling van \(\mathrm{Na^+}\) en \(\mathrm{Cl^-}\) bestaan.
\end{oplossing}

\begin{oefening}[Een XRD-verslag controleren: is dit silicium?]
Een student meet een poederdiffractogram van een grijs poeder uit een kapot zonnepaneel, met Cu K\(\alpha\)-straling. Het verslag vermeldt: ``eerste piek bij \(\theta = 28.44^\circ\), afkomstig van de \((111)\)-vlakken''. Voor silicium (diamantstructuur, 8 atomen per cel, \(M = 28.09\,\mathrm{g/mol}\)) is \(a = 0.5431\,\mathrm{nm}\) bekend.

\begin{enumerate}
\item[(a)]
  Bereken \(d_{111}\) en \(a\) wanneer je \(28.44^\circ\) letterlijk als \(\theta\) gebruikt.
\item[(b)]
  Herhaal in de veronderstelling dat het toestel \(2\theta = 28.44^\circ\) rapporteerde. Welke interpretatie is juist?
\item[(c)]
  Controleer je besluit met de dichtheid van silicium (\(2.33\,\mathrm{g/cm^3}\)).
\item[(d)]
  Bij welke \(2\theta\) verwacht je de \((220)\)-piek?
\end{enumerate}
\end{oefening}
\begin{oplossing}
\textbf{Resultaat:} (a) \(d = 0.162\,\mathrm{nm}\) en \(a = 0.281\,\mathrm{nm}\): onmogelijk. (b) Met \(\theta = 14.22^\circ\): \(d_{111} =\) \textbf{0.3136 nm} en \(a =\) \textbf{0.5431 nm}: dit is silicium; de gerapporteerde waarde was \(2\theta\). (c) Met \(a = 0.5431\,\mathrm{nm}\) volgt \textbf{2.33 g/cm\ensuremath{^{3}}}. (d) \(2\theta \approx\) \textbf{47.3°}.

\textbf{(a)} \(d = \dfrac{0.1541}{2\sin 28.44^\circ} = 0.162\,\mathrm{nm}\) en \(a = d\sqrt{3} = 0.281\,\mathrm{nm}\). Dat is bijna de helft van de bekende waarde, en de dichtheid zou dan acht keer te groot zijn.

\textbf{(b)} \(\theta = 14.22^\circ\): \(d_{111} = \dfrac{0.1541}{2\sin 14.22^\circ} = 0.3136\,\mathrm{nm}\) en \(a = 0.3136\sqrt{3} = 0.5431\,\mathrm{nm}\). Dit is exact de roosterconstante van silicium: het toestel rapporteerde \(2\theta\) en de student vergiste zich in het verslag (dezelfde valkuil als in het voorbeeld van §10.3.2).

\textbf{(c)} \(\rho = \dfrac{8 \times 28.09}{(6.022\times10^{23})(0.5431\times10^{-7})^3} = 2.33\,\mathrm{g/cm^3}\). De dichtheid is een snelle, onafhankelijke controle op een XRD-analyse.

\textbf{(d)} \(d_{220} = a/\sqrt{8} = 0.1920\,\mathrm{nm}\), \(\sin\theta = 0.1541/(2 \times 0.1920) = 0.401\), \(\theta = 23.66^\circ\), dus \(2\theta = 47.3^\circ\). Wie de reeks pieken met een databank vergelijkt, identificeert het materiaal in enkele seconden. Recyclagebedrijven gebruiken dit om silicium van hoge zuiverheid uit oude panelen terug te winnen.
\end{oplossing}

\begin{oefening}[Millerindices en het diffractogram van aluminium]
\begin{enumerate}
\item[(a)]
  Een vlak snijdt de kristalassen op \(x = a\), \(y = 2a\) en is evenwijdig met de \(z\)-as. Bepaal de Millerindices.
\item[(b)]
  Bereken voor aluminium (FCC, \(a = 0.405\,\mathrm{nm}\)) de vlakafstanden \(d_{100}\), \(d_{110}\), \(d_{111}\) en \(d_{200}\).
\item[(c)]
  Bereken de hoeken \(2\theta\) waaronder de \((111)\)- en \((200)\)-pieken verschijnen bij Cu K\(\alpha\)-straling.
\item[(d)]
  Welke vlakken liggen het dichtst bij elkaar? Leid af dat een hogere Millerindex een kleinere vlakafstand en een grotere diffractiehoek geeft.
\item[(e)]
  Je meet een aluminium blikje bij \(200\,^\circ\mathrm{C}\). Verschuiven de pieken naar grotere of kleinere \(2\theta\)?
\end{enumerate}
\end{oefening}
\begin{oplossing}
\textbf{Resultaat:} (a) \textbf{(210)}. (b) \(d_{100} = 0.405\), \(d_{110} = 0.286\), \(d_{111} =\) \textbf{0.234} en \(d_{200} = 0.203\,\mathrm{nm}\). (c) \((111)\) bij \textbf{38.5°}, \((200)\) bij \textbf{44.7°}. (d) \(d\) daalt met \(\sqrt{h^2 + k^2 + l^2}\). (e) Naar \textbf{kleinere} \(2\theta\) door thermische uitzetting.

\textbf{(a)} Snijpunten in eenheden van \(a\): \(1\), \(2\), \(\infty\). Omgekeerden: \(1\), \(\tfrac{1}{2}\), \(0\). Vermenigvuldigen met 2 tot de kleinste gehele getallen: \((210)\).

\textbf{(b)} Met \(d_{hkl} = a/\sqrt{h^2 + k^2 + l^2}\):

\begin{center}
\begin{tabular}{lll}
\toprule
\((hkl)\) & \(h^2 + k^2 + l^2\) & \(d\) (nm) \\
\midrule
\begin{enumerate}
 \def\labelenumi{(\arabic{enumi})} \setcounter{enumi}{99}  \item[(a)] \end{enumerate} & 1 & 0.405 \\
\begin{enumerate}
 \def\labelenumi{(\arabic{enumi})} \setcounter{enumi}{109}  \item[(a)] \end{enumerate} & 2 & 0.286 \\
\begin{enumerate}
 \def\labelenumi{(\arabic{enumi})} \setcounter{enumi}{110}  \item[(a)] \end{enumerate} & 3 & 0.234 \\
\begin{enumerate}
 \def\labelenumi{(\arabic{enumi})} \setcounter{enumi}{199}  \item[(a)] \end{enumerate} & 4 & 0.203 \\
\bottomrule
\end{tabular}
\end{center}

\textbf{(c)} \((111)\): \(\sin\theta = 0.1541/(2 \times 0.2338) = 0.330\), \(\theta = 19.25^\circ\), \(2\theta = 38.5^\circ\). \((200)\): \(\sin\theta = 0.1541/(2 \times 0.2025) = 0.381\), \(2\theta = 44.7^\circ\). Dit zijn inderdaad de twee sterkste pieken in een diffractogram van aluminium. (Bij een FCC-rooster heffen de reflecties van de \((100)\)- en \((110)\)-vlakken elkaar uit door de extra atomen op de vlakken; daarom verschijnen ze niet.)

\textbf{(d)} Hoe groter \(h^2 + k^2 + l^2\), hoe kleiner \(d\), en volgens Bragg hoe groter \(\sin\theta = \lambda/(2d)\). Vlakken met hoge indices liggen dicht op elkaar, maar bevatten ook minder atomen per oppervlakte.

\textbf{(e)} Bij opwarmen zet het rooster uit: \(a\) en alle \(d\) worden groter, dus \(\sin\theta\) kleiner. De pieken schuiven naar kleinere hoeken. Met hoge-temperatuur-XRD meet men zo de thermische uitzettingscoëfficiënt van een materiaal, of volgt men een faseovergang zoals die van oefening 10.2.
\end{oplossing}

\begin{oefening}[Vacatures in aluminium: waarom gloeien werkt]
De vormingsenergie van een vacature in aluminium is \(E_v = 0.76\,\mathrm{eV}\); het smeltpunt is \(660\,^\circ\mathrm{C}\). De evenwichtsfractie vacatures is \(n_v/N = \exp(-E_v/k_BT)\).

\begin{enumerate}
\item[(a)]
  Bereken de vacaturefractie bij \(300\,\mathrm{K}\), \(600\,\mathrm{K}\) en \(900\,\mathrm{K}\).
\item[(b)]
  Hoeveel vacatures bevat \(1\,\mathrm{cm^3}\) aluminium (\(4\) atomen per cel van \((0.405\,\mathrm{nm})^3\)) bij \(300\) en bij \(900\,\mathrm{K}\)?
\item[(c)]
  Atomen diffunderen in een vaste stof door in een naburige vacature te springen. Verklaar waarom de diffusiesnelheid in metalen nog sterker stijgt met de temperatuur dan de vacaturefractie alleen. Verbind met de vergelijking van Arrhenius (hoofdstuk 8).
\item[(d)]
  Een plaatwerker gloeit koudvervormd, hard geworden aluminium enkele minuten bij \(350\,^\circ\mathrm{C}\) om het weer zacht te maken. Waarom gebeurt dat niet vanzelf bij kamertemperatuur?
\end{enumerate}
\end{oefening}
\begin{oplossing}
\textbf{Resultaat:} (a) \textbf{\(1.7\times10^{-13}\)}, \textbf{\(4.1\times10^{-7}\)} en \textbf{\(5.5\times10^{-5}\)}. (b) Ongeveer \textbf{\(10^{10}\)} tegenover \textbf{\(3.3\times10^{18}\)} vacatures per cm\ensuremath{^{3}}. (c) Zowel het aantal vacatures als de springfrequentie volgen een Arrheniusterm. (d) Bij kamertemperatuur zijn atomen \textbf{praktisch onbeweeglijk}.

\textbf{(a)} \(k_BT\) is \(0.0259\,\mathrm{eV}\) bij \(300\,\mathrm{K}\), \(0.0517\,\mathrm{eV}\) bij \(600\,\mathrm{K}\) en \(0.0776\,\mathrm{eV}\) bij \(900\,\mathrm{K}\):

\begin{itemize}
\item
  \(300\,\mathrm{K}\): \(e^{-29.4} = 1.7\times10^{-13}\)
\item
  \(600\,\mathrm{K}\): \(e^{-14.7} = 4.1\times10^{-7}\)
\item
  \(900\,\mathrm{K}\): \(e^{-9.8} = 5.5\times10^{-5}\)
\end{itemize}

Tussen kamertemperatuur en net onder het smeltpunt stijgt de fractie met een factor \(3\times10^{8}\).

\textbf{(b)} \(N = 4/(0.405\times10^{-7}\,\mathrm{cm})^3 = 6.0\times10^{22}\) atomen per cm\ensuremath{^{3}}. Bij \(300\,\mathrm{K}\): \(1.0\times10^{10}\) vacatures; bij \(900\,\mathrm{K}\): \(3.3\times10^{18}\). Gemiddeld is dan één op de \(18000\) roosterplaatsen leeg.

\textbf{(c)} Een atoom kan enkel springen als er een vacature naast zit (factor \(e^{-E_v/k_BT}\)) én als het de energiebarrière voor de sprong zelf overwint (een tweede factor \(e^{-E_m/k_BT}\), met \(E_m\) de migratie-energie, voor Al ongeveer \(0.6\,\mathrm{eV}\)). De diffusiecoëfficiënt volgt dus een Arrheniusvergelijking met een schijnbare activeringsenergie \(E_v + E_m \approx 1.4\,\mathrm{eV}\), ongeveer \(130\,\mathrm{kJ/mol}\): een uitgesproken temperatuurgevoelig proces, net zoals een chemische reactie met een hoge \(E_a\).

\textbf{(d)} Koudvervormen brengt veel dislocaties in het metaal; die hinderen elkaar en maken het hard (§10.4.2). Om ze te laten verdwijnen moeten atomen zich verplaatsen, en dat vraagt diffusie. Bij kamertemperatuur is diffusie in aluminium zo traag dat er in jaren niets merkbaars gebeurt; bij \(350\,^\circ\mathrm{C}\) gaat het een factor \(10^{10}\) of meer sneller en vormt zich in minuten een nieuwe, zachte korrelstructuur (rekristallisatie).
\end{oplossing}

\begin{oefening}[Wustiet: ijzeroxide dat niet klopt met zijn formule]
Op gloeiend staal vormt zich een dikke oxidelaag (walshuid) die grotendeels uit wustiet bestaat, met de samenstelling \(\mathrm{Fe_{0.95}O}\). Het kristal heeft de NaCl-structuur, waarin een deel van de ijzerplaatsen leeg is.

\begin{enumerate}
\item[(a)]
  Welke fractie van de ijzerplaatsen is leeg?
\item[(b)]
  De zuurstofionen zijn \(\mathrm{O^{2-}}\). Welk deel van de ijzerionen moet \(\mathrm{Fe^{3+}}\) zijn om het kristal elektrisch neutraal te houden?
\item[(c)]
  Waarom geleidt wustiet elektriciteit, terwijl een perfect stoichiometrisch ionair kristal een isolator zou zijn?
\item[(d)]
  Leg uit waarom zo'n defecte oxidelaag het staal niet beschermt tegen verdere oxidatie, terwijl de \(\mathrm{Cr_2O_3}\)-film op roestvast staal (§10.8) dat wel doet.
\end{enumerate}
\end{oefening}
\begin{oplossing}
\textbf{Resultaat:} (a) \textbf{5 \%} van de ijzerplaatsen is leeg. (b) \(0.10\) van de \(0.95\) ijzerionen, dus ongeveer \textbf{10.5 \%}, is \(\mathrm{Fe^{3+}}\). (c) Elektronen springen tussen \(\mathrm{Fe^{2+}}\) en \(\mathrm{Fe^{3+}}\). (d) Via de \textbf{vacatures} diffundeert ijzer snel naar buiten.

\textbf{(a)} Per zuurstofion is er \(0.95\) ijzer in plaats van \(1\): \(5\,\%\) van de kationplaatsen is een vacature.

\textbf{(b)} Noem \(y\) het aantal \(\mathrm{Fe^{3+}}\) per formule-eenheid. De lading moet \(+2\) zijn: \(3y + 2(0.95 - y) = 2\), dus \(y = 0.10\). Van de \(0.95\) ijzerionen is dus \(0.10/0.95 = 10.5\,\%\) driewaardig. Elke vacature wordt gecompenseerd door twee \(\mathrm{Fe^{3+}}\)-ionen.

\textbf{(c)} Een elektron kan van een \(\mathrm{Fe^{2+}}\) naar een naburig \(\mathrm{Fe^{3+}}\) springen, waardoor beide van rol wisselen: het ``gat'' beweegt door het kristal. Wustiet is daardoor een p-type halfgeleider, net als silicium gedoteerd met boor (§10.5.2), maar hier zijn het de vacatures zelf die de ladingsdragers leveren.

\textbf{(d)} Door de vele vacatures kunnen \(\mathrm{Fe^{2+}}\)-ionen gemakkelijk van het metaal door de oxidelaag naar buiten diffunderen, waar ze met zuurstof reageren. De laag groeit dus snel verder. Chroomoxide is een veel stoichiometrischer en dichter oxide met weinig defecten: ionen en elektronen geraken er nauwelijks door, zodat de film na enkele nanometer stopt met groeien en het staal passief blijft (hoofdstuk 7).
\end{oplossing}

\begin{oefening}[Welke kleur licht? Band gap, zonnecel en led]
Een foton kan een elektron van de valentieband naar de geleidingsband brengen als zijn energie minstens gelijk is aan \(E_g\). Gebruik \(E = hc/\lambda\) (hoofdstuk 2).

\begin{enumerate}
\item[(a)]
  Bereken de maximale golflengte die nog geabsorbeerd wordt door Si (\(1.12\,\mathrm{eV}\)), GaAs (\(1.42\,\mathrm{eV}\)), \(\mathrm{TiO_2}\) (\(3.2\,\mathrm{eV}\)) en diamant (\(5.5\,\mathrm{eV}\)).
\item[(b)]
  Zichtbaar licht ligt tussen \(400\) en \(700\,\mathrm{nm}\). Verklaar waarom silicium donkergrijs is, diamant doorzichtig en \(\mathrm{TiO_2}\) wit.
\item[(c)]
  Waarom zit \(\mathrm{TiO_2}\) in zonnecrème en in zelfreinigend glas?
\item[(d)]
  Een led zendt licht uit wanneer een elektron over de band gap terugvalt. Welke band gap is nodig voor een rode led (\(630\,\mathrm{nm}\)) en voor een blauwe led (\(450\,\mathrm{nm}\))? Waarom was de blauwe led zo moeilijk te maken (Nobelprijs 2014)?
\end{enumerate}
\end{oefening}
\begin{oplossing}
\textbf{Resultaat:} (a) Si \textbf{1107 nm}, GaAs \textbf{873 nm}, \(\mathrm{TiO_2}\) \textbf{388 nm}, diamant \textbf{225 nm}. (b) Si absorbeert al het zichtbare licht, diamant niets, \(\mathrm{TiO_2}\) verstrooit alles. (c) \(\mathrm{TiO_2}\) absorbeert \textbf{UV}. (d) Rood: \textbf{1.97 eV}; blauw: \textbf{2.76 eV}, wat een moeilijk te kweken materiaal (GaN) vereiste.

\textbf{(a)} \(\lambda_{\max} = hc/E_g = 1240/E_g\) (nm, met \(E_g\) in eV):

\begin{itemize}
\item
  Si: \(1240/1.12 = 1107\,\mathrm{nm}\) (nabij-infrarood)
\item
  GaAs: \(873\,\mathrm{nm}\)
\item
  \(\mathrm{TiO_2}\): \(388\,\mathrm{nm}\) (UV)
\item
  diamant: \(225\,\mathrm{nm}\) (ver UV)
\end{itemize}

\textbf{(b)} Silicium absorbeert alle zichtbare fotonen (allemaal energierijker dan \(1.12\,\mathrm{eV}\)) en ziet er daarom donker, metaalachtig grijs uit. Diamant absorbeert geen enkel zichtbaar foton en is doorzichtig. \(\mathrm{TiO_2}\) absorbeert evenmin zichtbaar licht, maar heeft een zeer hoge brekingsindex: een poeder van fijne korrels verstrooit alle kleuren gelijk en ziet er wit uit. Het is het witte pigment in muurverf, tandpasta en papier.

\textbf{(c)} In zonnecrème absorbeert \(\mathrm{TiO_2}\) de UV-straling onder \(388\,\mathrm{nm}\) die huidkanker veroorzaakt, terwijl het zichtbare licht doorgaat. Op glas gebruikt men de geabsorbeerde UV-energie: het elektron in de geleidingsband en het achtergebleven gat zijn sterke reductoren en oxidatoren, die organisch vuil afbreken (fotokatalyse, hoofdstuk 8). Onder zichtbaar licht werkt dit niet: de fotonen hebben te weinig energie.

\textbf{(d)} \(E_g = 1240/630 = 1.97\,\mathrm{eV}\) voor rood en \(1240/450 = 2.76\,\mathrm{eV}\) voor blauw. Rode leds (op basis van GaAsP) bestaan sinds de jaren zestig. Voor blauw is een halfgeleider met een grote band gap nodig, zoals GaN (\(3.4\,\mathrm{eV}\), met indium bijgemengd om de gap te verkleinen). Het lukte decennia niet om GaN-kristallen met weinig defecten te kweken en om er p-type doping in te krijgen. Dat dit uiteindelijk wel lukte, maakte witte ledlampen mogelijk.
\end{oplossing}

\begin{oefening}[Hoeveel fosfor maakt van silicium een geleider?]
Zuiver silicium (\(a = 0.5431\,\mathrm{nm}\), 8 atomen per cel) bevat bij \(300\,\mathrm{K}\) ongeveer \(1\times10^{10}\) vrije elektronen per cm\ensuremath{^{3}}. Een chipfabrikant doteert het met \(1\,\mathrm{ppm}\) fosfor (één P-atoom per miljoen Si-atomen).

\begin{enumerate}
\item[(a)]
  Is dit n- of p-type doping? Verklaar met het aantal valentie-elektronen.
\item[(b)]
  Bereken het aantal Si-atomen per cm\ensuremath{^{3}} en de elektronenconcentratie na doping, als elk P-atoom een elektron afstaat. Met welke factor stijgt het aantal ladingsdragers?
\item[(c)]
  Vergelijk de Boltzmannfactor \(e^{-E/k_BT}\) bij \(300\,\mathrm{K}\) voor het overbruggen van de band gap (\(1.12\,\mathrm{eV}\)) en voor het ioniseren van een P-donor (\(0.045\,\mathrm{eV}\)). Waarom zijn bijna alle donoren geïoniseerd?
\item[(d)]
  Waarom moet silicium voor chips eerst tot een zuiverheid van ``elf negens'' (99.999999999 \%) gezuiverd worden?
\end{enumerate}
\end{oefening}
\begin{oplossing}
\textbf{Resultaat:} (a) \textbf{n-type}: P heeft 5 valentie-elektronen. (b) \(5.0\times10^{22}\) Si-atomen per cm\ensuremath{^{3}}; \textbf{\(5.0\times10^{16}\)}** elektronen per cm\ensuremath{^{3}}\textbf{, een factor }\(5\times10^{6}\)** meer. (c) \(1.5\times10^{-19}\) tegenover \textbf{0.18}. (d) Ongewenste onzuiverheden zouden zelf als dopering werken.

\textbf{(a)} Silicium heeft vier valentie-elektronen, die alle vier in bindingen met de buren zitten. Een P-atoom op een Si-plaats heeft er vijf; het vijfde is zwak gebonden en komt gemakkelijk in de geleidingsband: n-type.

\textbf{(b)} \(n_{\mathrm{Si}} = 8/(0.5431\times10^{-7}\,\mathrm{cm})^3 = 5.0\times10^{22}\,\mathrm{cm^{-3}}\). Met \(1\,\mathrm{ppm}\): \(5.0\times10^{16}\) donoren, en dus evenveel vrije elektronen per cm\ensuremath{^{3}}. Dat is vijf miljoen keer meer dan in zuiver silicium, en de geleidbaarheid stijgt evenredig.

\textbf{(c)} \(k_BT = 0.0259\,\mathrm{eV}\):

\begin{itemize}
\item
  band gap: \(e^{-1.12/0.0259} = 1.5\times10^{-19}\)
\item
  donor: \(e^{-0.045/0.0259} = 0.18\)
\end{itemize}

De donorenergie is van dezelfde grootte als \(k_BT\): thermische trillingen ioniseren vrijwel alle donoren. Over de volledige band gap springt bijna geen elektron. Daarom bepaalt de doping, en niet de temperatuur, het aantal ladingsdragers in een chip bij kamertemperatuur.

\textbf{(d)} Elke vreemde stof die \(1\,\mathrm{ppm}\) of zelfs \(1\,\mathrm{ppb}\) bedraagt, levert al veel meer ladingsdragers dan zuiver silicium zelf. Om de geleiding nauwkeurig te kunnen instellen via gecontroleerde doping, moeten ongecontroleerde onzuiverheden ver onder dat niveau liggen. Dat zuiveren gebeurt door destillatie van \(\mathrm{SiHCl_3}\) en door zonesmelten: onzuiverheden lossen beter op in de smelt dan in het kristal (een fasediagram zoals in hoofdstuk 5).
\end{oplossing}

\begin{oefening}[Gietijzer of aluminium? Warmtecapaciteit van een kookpot]
\begin{enumerate}
\item[(a)]
  Schat met de regel van Dulong--Petit de specifieke warmtecapaciteit (in \(\mathrm{J/(g\cdot K)}\)) van ijzer, aluminium, koper en lood. Vergelijk met de tabelwaarden \(0.449\), \(0.897\), \(0.385\) en \(0.129\,\mathrm{J/(g\cdot K)}\).
\item[(b)]
  Hoeveel warmte is nodig om een gietijzeren pan van \(2.0\,\mathrm{kg}\) en een aluminium pan van \(1.0\,\mathrm{kg}\) van \(20\) tot \(200\,^\circ\mathrm{C}\) op te warmen?
\item[(c)]
  Waarom heeft een metaal met zware atomen per gram een kleine warmtecapaciteit?
\item[(d)]
  Waarom wijkt diamant sterk af van Dulong--Petit bij kamertemperatuur, maar steeds minder bij hogere temperatuur?
\end{enumerate}
\end{oefening}
\begin{oplossing}
\textbf{Resultaat:} (a) Fe \textbf{0.447}, Al \textbf{0.924}, Cu \textbf{0.392}, Pb \textbf{0.120 J/(g\ensuremath{\cdot{}}K)}, binnen enkele procenten van de tabelwaarden. (b) Beide ongeveer \textbf{160--165 kJ}. (c) \(c = 3R/M\): per mol is \(C\) gelijk, dus per gram omgekeerd evenredig met \(M\). (d) Trillingen met een hoge frequentie worden pas bij hoge temperatuur thermisch actief.

\textbf{(a)} \(c = 3R/M\):

\begin{center}
\begin{tabular}{llll}
\toprule
Metaal & \(M\) (g/mol) & \(3R/M\) & tabel \\
\midrule
Fe & 55.85 & 0.447 & 0.449 \\
Al & 26.98 & 0.924 & 0.897 \\
Cu & 63.55 & 0.392 & 0.385 \\
Pb & 207.2 & 0.120 & 0.129 \\
\bottomrule
\end{tabular}
\end{center}

De afwijkingen zijn klein; voor lood ligt de echte waarde hoger, omdat \(C_p\) iets groter is dan \(C_V\) en lood bij kamertemperatuur al dicht bij zijn smeltpunt zit.

\textbf{(b)} Gietijzer: \(Q = (2000\,\mathrm{g})(0.447)(180\,\mathrm{K}) = 161\,\mathrm{kJ}\). Aluminium: \(Q = (1000)(0.924)(180) = 166\,\mathrm{kJ}\). De twee pannen vragen vrijwel evenveel energie: de zware gietijzeren pan bevat ongeveer evenveel mol atomen als de lichte aluminium pan (\(36\) tegenover \(37\,\mathrm{mol}\)). Gietijzer houdt de warmte wel langer vast bij het dichtschroeien van vlees, omdat de warmte minder snel naar de rand weglekt (lagere warmtegeleiding).

\textbf{(c)} Elk atoom slaat volgens het equipartitiebeginsel dezelfde energie op, \(3k_BT\), ongeacht zijn massa. Per gram zijn er minder zware atomen, dus is er minder energie nodig per graad. Daarom warmt een loden gewicht in je hand zo snel op.

\textbf{(d)} De sterke, lichte C--C-bindingen trillen met hoge frequenties. Een trilling neemt pas volledig deel aan de warmteopslag wanneer \(k_BT\) groot is ten opzichte van de energiekwanta van die trilling (\(h\nu\), hoofdstuk 2). Bij kamertemperatuur is dat voor diamant niet het geval. Bij een temperatuur van meer dan \(2000\,\mathrm{K}\) nadert ook diamant de waarde \(3R\).
\end{oplossing}

\begin{oefening}[Hoeveel energie bespaart het recycleren van een blikje?]
\begin{enumerate}
\item[(a)]
  Schat met de regel van Richards de smeltenthalpie van aluminium (\(T_{\text{fus}} = 933\,\mathrm{K}\)), koper (\(1358\,\mathrm{K}\)), lood (\(601\,\mathrm{K}\)) en ijzer (\(1811\,\mathrm{K}\)). Vergelijk met de tabelwaarden \(10.7\), \(13.3\), \(4.8\) en \(13.8\,\mathrm{kJ/mol}\).
\item[(b)]
  Bereken de minimale energie om \(1\,\mathrm{kg}\) aluminium blikjes van \(25\,^\circ\mathrm{C}\) op te warmen en te smelten. Gebruik Dulong--Petit en de tabelwaarde voor de smeltenthalpie. Geef het resultaat in kWh.
\item[(c)]
  Primair aluminium wordt elektrolytisch uit aluminiumoxide gemaakt (proces van Hall--Héroult), met een verbruik van ongeveer \(14\,\mathrm{kWh}\) per kg. Welk percentage energie bespaart recyclage, in theorie?
\item[(d)]
  Waarom verbruikt een smeltoven in de praktijk wel twee à drie keer meer dan je berekende waarde?
\end{enumerate}
\end{oefening}
\begin{oplossing}
\textbf{Resultaat:} (a) Al \textbf{7.8}, Cu \textbf{11.4}, Pb \textbf{5.0}, Fe \textbf{15.2 kJ/mol}: binnen ongeveer 10 tot 25 \%. (b) Ongeveer \textbf{0.98 MJ = 0.27 kWh} per kg. (c) Meer dan \textbf{98 \%} besparing. (d) Warmteverliezen via rookgassen, wanden en het verwarmen van de oven zelf.

\textbf{(a)} \(\Delta H_{\text{fus}} \approx 8.4\,T_{\text{fus}}\) (in J/mol):

\begin{center}
\begin{tabular}{llll}
\toprule
Metaal & \(T_{\text{fus}}\) (K) & Richards (kJ/mol) & tabel (kJ/mol) \\
\midrule
Al & 933 & 7.8 & 10.7 \\
Cu & 1358 & 11.4 & 13.3 \\
Pb & 601 & 5.0 & 4.8 \\
Fe & 1811 & 15.2 & 13.8 \\
\bottomrule
\end{tabular}
\end{center}

Een eenvoudige vuistregel die de grootteorde goed weergeeft, maar per metaal tot \(25\,\%\) kan afwijken.

\textbf{(b)} Opwarmen tot het smeltpunt: \((1000\,\mathrm{g})(0.924\,\mathrm{J/(g\cdot K)})(933 - 298\,\mathrm{K}) = 587\,\mathrm{kJ}\). Smelten: \((1000/26.98\,\mathrm{mol})(10.7\,\mathrm{kJ/mol}) = 397\,\mathrm{kJ}\). Totaal \(984\,\mathrm{kJ} = 0.27\,\mathrm{kWh}\), ongeveer de energie om een volle waterkoker twee keer aan de kook te brengen.

\textbf{(c)} \(1 - 0.27/14 = 0.98\): recyclage vraagt in theorie minder dan \(2\,\%\) van de energie van primair aluminium. Bij de elektrolyse moet de sterke Al--O-binding verbroken worden (een grote, positieve \(\Delta_rG^\circ\), hoofdstuk 7); bij het hersmelten blijft het metaal metaal en is enkel de warmte voor opwarmen en smelten nodig.

\textbf{(d)} De oven, de bekleding en de rookgassen moeten mee opgewarmd worden, het gesmolten aluminium moet boven het smeltpunt gehouden worden om te kunnen gieten, en lak en plastic op de blikjes moeten afgebrand worden. Realistisch verbruikt hersmelten zo'n \(0.5\)--\(0.8\,\mathrm{kWh/kg}\), nog altijd ongeveer \(95\,\%\) minder dan primair aluminium.
\end{oplossing}

\begin{oefening}[Geïntegreerde case: welke kathode voor een batterij?]
Twee veelgebruikte kathodematerialen voor Li-ionbatterijen zijn het gelaagde \(\mathrm{LiCoO_2}\) (\(M = 97.87\,\mathrm{g/mol}\), gemiddelde spanning \(3.9\,\mathrm{V}\)) en het olivijn \(\mathrm{LiFePO_4}\) (\(M = 157.76\,\mathrm{g/mol}\), \(3.4\,\mathrm{V}\)). Bij het laden verlaat lithium het rooster: \(\mathrm{LiCoO_2 \rightarrow Li_{1-x}CoO_2 + x\,Li^+ + x\,e^-}\).

\begin{enumerate}
\item[(a)]
  Bereken de theoretische capaciteit (in mAh/g) van beide materialen als alle lithium uitgewisseld kan worden.
\item[(b)]
  Bij \(\mathrm{LiCoO_2}\) kan in de praktijk slechts ongeveer de helft van het lithium verwijderd worden (\(x \approx 0.5\)), bij \(\mathrm{LiFePO_4}\) vrijwel alles. Bereken de bruikbare capaciteit en de specifieke energie (Wh per kg kathodemateriaal).
\item[(c)]
  Hoeveel gram kathodemateriaal bevat een smartphonebatterij van \(15.4\,\mathrm{Wh}\) (oefening 7.4) minstens, voor elk materiaal?
\item[(d)]
  Waarom kan men bij \(\mathrm{LiCoO_2}\) niet al het lithium verwijderen? Gebruik de begrippen non-stoichiometrie en kristalstructuur.
\item[(e)]
  Waarom kiezen fabrikanten van elektrische bussen en thuisbatterijen vaak voor \(\mathrm{LiFePO_4}\), terwijl smartphones \(\mathrm{LiCoO_2}\)-achtige materialen gebruiken?
\end{enumerate}
\end{oefening}
\begin{oplossing}
\textbf{Resultaat:} (a) \textbf{274 mAh/g} (\(\mathrm{LiCoO_2}\)) en \textbf{170 mAh/g} (\(\mathrm{LiFePO_4}\)). (b) Bruikbaar \(137\) en \(170\,\mathrm{mAh/g}\); \textbf{534} en \textbf{578 Wh/kg}. (c) Ongeveer \textbf{29 g} en \textbf{27 g}. (d) Een te sterk ontladen rooster wordt \textbf{onstabiel}. (e) \(\mathrm{LiFePO_4}\) is \textbf{veiliger, goedkoper en gaat langer mee}; \(\mathrm{LiCoO_2}\) is compacter.

\textbf{(a)} Eén elektron per formule-eenheid: \(Q = F/M\), en \(1\,\mathrm{mAh} = 3.6\,\mathrm{C}\).

\begin{itemize}
\item
  \(\mathrm{LiCoO_2}\): \(96485/97.87 = 986\,\mathrm{C/g} = 274\,\mathrm{mAh/g}\)
\item
  \(\mathrm{LiFePO_4}\): \(96485/157.76 = 612\,\mathrm{C/g} = 170\,\mathrm{mAh/g}\)
\end{itemize}

\textbf{(b)} \(\mathrm{LiCoO_2}\): \(0.5 \times 274 = 137\,\mathrm{mAh/g}\), maal \(3.9\,\mathrm{V}\) geeft \(534\,\mathrm{Wh/kg}\). \(\mathrm{LiFePO_4}\): \(170 \times 3.4 = 578\,\mathrm{Wh/kg}\). Per kg actief materiaal doen ze het dus ongeveer even goed; de hogere spanning van kobalt compenseert het kleinere bruikbare aandeel lithium.

\textbf{(c)} \(15.4/0.534 = 29\,\mathrm{g}\) \(\mathrm{LiCoO_2}\), of \(15.4/0.578 = 27\,\mathrm{g}\) \(\mathrm{LiFePO_4}\). Een echte cel is zwaarder: anode, elektrolyt, separator, stroomcollectoren en behuizing wegen samen meer dan de kathode.

\textbf{(d)} \(\mathrm{Li_{1-x}CoO_2}\) is een non-stoichiometrische verbinding (§10.4.4): tussen de \(\mathrm{CoO_2}\)-lagen zitten lithiumionen die de lagen elektrostatisch bijeenhouden. Haal je meer dan ongeveer de helft weg, dan stoten de negatief geladen zuurstoflagen elkaar af, verschuiven de lagen en verandert de kristalstructuur. Bovendien wordt \(\mathrm{Co^{4+}}\) zo'n sterke oxidator dat het zuurstof uit het rooster vrijzet en de elektrolyt aantast. Dat beschadigt de batterij onomkeerbaar en kan bij oververhitting tot brand leiden.

\textbf{(e)} In het olivijnrooster van \(\mathrm{LiFePO_4}\) zit zuurstof stevig gebonden in fosfaatgroepen; ook volledig ontladen (\(\mathrm{FePO_4}\)) blijft de structuur stabiel en komt er geen zuurstof vrij. Het materiaal is daardoor veiliger, gaat duizenden cycli mee en bevat geen duur en schaars kobalt. Nadelen: een lagere spanning en een lagere dichtheid (\(3.6\) tegenover \(5.1\,\mathrm{g/cm^3}\)), zodat een \(\mathrm{LiFePO_4}\)-batterij voor dezelfde energie meer plaats inneemt. In een smartphone telt elke kubieke millimeter; in een bus of een thuisbatterij tellen veiligheid, levensduur en prijs meer.
\end{oplossing}

\begin{oefening}[Tinpest: wanneer valt tin uiteen?]
Wit tin (\ensuremath{\beta{}}-tin, \(\rho = 7.29\,\mathrm{g/cm^3}\)) is de gewone metallische vorm. Onder \(13.2\,^\circ\mathrm{C}\) is grijs tin (\ensuremath{\alpha{}}-tin, diamantstructuur, \(\rho = 5.77\,\mathrm{g/cm^3}\)) de stabiele vorm. Voor de overgang wit \ensuremath{\rightarrow{}} grijs bedraagt \(\Delta_{\text{trs}}H = -2.09\,\mathrm{kJ/mol}\). Gebruik \(M = 118.71\,\mathrm{g/mol}\). Volgens een bekend verhaal vielen de tinnen knopen van Napoleons soldaten in de Russische winter van 1812 uit elkaar.

\begin{enumerate}
\item[(a)]
  Bereken \(\Delta_{\text{trs}}S\) voor de overgang wit \ensuremath{\rightarrow{}} grijs. Verklaar het teken.
\item[(b)]
  Bereken \(\Delta_{\text{trs}}G\) bij \(-30\,^\circ\mathrm{C}\). Is de omzetting spontaan?
\item[(c)]
  Bereken de relatieve volumetoename bij de overgang. Waarom valt een voorwerp van tin daardoor tot poeder uiteen?
\item[(d)]
  Bereken met de vergelijking van Clapeyron hoeveel de overgangstemperatuur verandert per bar extra druk.
\item[(e)]
  Wit tin kan in de praktijk jaren onder \(13\,^\circ\mathrm{C}\) bestaan. Waarom? Wat betekent dat voor het verhaal over Napoleon?
\end{enumerate}
\end{oefening}
\begin{oplossing}
\textbf{Resultaat:} (a) \(\Delta_{\text{trs}}S\) = \textbf{\ensuremath{-}7.3 J/(mol\ensuremath{\cdot{}}K)}. (b) \(\Delta_{\text{trs}}G\) = \textbf{\ensuremath{-}0.32 kJ/mol}: spontaan. (c) \textbf{+26 \%} volume. (d) \(\mathrm{d}p/\mathrm{d}T\) = \ensuremath{-}17 bar/K, dus \textbf{\ensuremath{-}0.06 K per bar}. (e) De omzetting is kinetisch zeer traag. Het verhaal is waarschijnlijk overdreven.

\textbf{(a)} Bij de overgangstemperatuur is \(\Delta G = 0\), dus \(\Delta_{\text{trs}}S = \Delta_{\text{trs}}H/T_{\text{trs}} = -2090/286.35 = -7.3\,\mathrm{J/(mol\cdot K)}\). Grijs tin heeft een stijve, covalente diamantstructuur waarin de atomen minder vrij kunnen trillen dan in het metallische witte tin, dus een lagere entropie.

\textbf{(b)} \(\Delta_{\text{trs}}G = \Delta H - T\Delta S = -2090 - (243.15)(-7.3) = -315\,\mathrm{J/mol}\). De omzetting is spontaan, maar de drijvende kracht is klein.

\textbf{(c)} \(V_m(\text{wit}) = 118.71/7.29 = 16.28\,\mathrm{cm^3/mol}\) en \(V_m(\text{grijs}) = 118.71/5.77 = 20.57\,\mathrm{cm^3/mol}\): een toename met 26 \%. De omzetting start op enkele kiemen en groeit vandaar naar buiten. Het uitzettende grijze tin legt enorme spanningen op het omliggende witte tin. Grijs tin is bovendien bros (covalent, zoals silicium), zodat het voorwerp verbrokkelt.

\textbf{(d)}

\[
\frac{\mathrm{d}p}{\mathrm{d}T} = \frac{\Delta_{\text{trs}}H}{T\,\Delta V} = \frac{-2090}{(286.35)(4.29\times10^{-6})} = -1.7\times10^{6}\,\mathrm{Pa/K} = -17\,\mathrm{bar/K}
\]

Per bar extra druk daalt de overgangstemperatuur met \(1/17 = 0.06\,\mathrm{K}\). Druk bevoordeelt de dichtere witte vorm.

\textbf{(e)} De vorming van een kiem van een nieuwe kristalstructuur vraagt een grote activeringsenergie. Bij lage temperatuur verloopt de groei bovendien traag (hoofdstuk 8). Thermodynamisch spontaan betekent dus niet snel. Zuiver tin moet vaak maanden onder \(-30\,^\circ\mathrm{C}\) blijven voordat tinpest zichtbaar wordt, en kleine hoeveelheden lood of antimoon onderdrukken ze. Of de knopen werkelijk een rol speelden in de nederlaag, is historisch niet aangetoond. Het verschijnsel is wel echt: het teisterde orgelpijpen in koude kerken en is een aandachtspunt bij loodvrij soldeer.
\end{oplossing}

\begin{oefening}[Hoe werkt een digitale koortsthermometer?]
Veel thermometers gebruiken een NTC-thermistor: een halfgeleider waarvan de weerstand daalt als de temperatuur stijgt. Voor een intrinsieke halfgeleider is de geleidbaarheid bij benadering evenredig met \(\exp(-E_g/2k_BT)\).

\begin{enumerate}
\item[(a)]
  Met welke factor neemt de geleidbaarheid van silicium (\(E_g = 1.12\,\mathrm{eV}\)) toe tussen \(25\) en \(100\,^\circ\mathrm{C}\)?
\item[(b)]
  Vergelijk de geleidbaarheid van silicium en diamant (\(E_g = 5.5\,\mathrm{eV}\)) bij \(25\,^\circ\mathrm{C}\), met dezelfde voorfactor. Hoeveel ordes van grootte verschillen ze?
\item[(c)]
  Hoe verandert de weerstand van een metaal zoals koper bij opwarming? Waarom is dat gedrag omgekeerd?
\item[(d)]
  Waarom is een sterk gedoteerde halfgeleider minder geschikt als thermistor dan een intrinsieke?
\end{enumerate}
\end{oefening}
\begin{oplossing}
\textbf{Resultaat:} (a) Ongeveer \textbf{80 keer} groter. (b) Silicium geleidt ongeveer \textbf{10\ensuremath{^{37}} keer} beter. (c) De weerstand van koper stijgt licht. (d) Bij doping hangt het aantal ladingsdragers nauwelijks af van de temperatuur.

\textbf{(a)} Met \(k_B = 8.617\times10^{-5}\,\mathrm{eV/K}\):

\[
\frac{\sigma(373\,\mathrm{K})}{\sigma(298\,\mathrm{K})} = \exp\left[-\frac{1.12}{2(8.617\times10^{-5})}\left(\frac{1}{373.15}-\frac{1}{298.15}\right)\right] = \exp(4.38) \approx 80
\]

Elke graad Celsius verandert de weerstand met enkele procenten: daardoor kan een thermistor temperaturen tot op \(0.1\,\mathrm{K}\) nauwkeurig meten.

\textbf{(b)} \(\sigma_{\mathrm{Si}}/\sigma_{\text{diamant}} = \exp[(5.5-1.12)/(2 \times 8.617\times10^{-5}\times298.15)] = \exp(85) \approx 10^{37}\). Een verschil in band gap van een factor vijf maakt van silicium een halfgeleider en van diamant een uitstekende isolator.

\textbf{(c)} In een metaal is de geleidingsband al deels gevuld (\(E_g = 0\)), dus het aantal ladingsdragers verandert niet met de temperatuur. Bij hogere temperatuur trillen de atomen sterker en verstrooien ze de elektronen vaker, zodat de weerstand licht stijgt (PTC-gedrag). Bij een halfgeleider overheerst de exponentiële toename van het aantal ladingsdragers.

\textbf{(d)} De donor- of acceptorniveaus liggen zo dicht bij de banden dat ze bij kamertemperatuur al bijna volledig geïoniseerd zijn. Het aantal ladingsdragers ligt dan vast door de dopingconcentratie en hangt nauwelijks af van de temperatuur. Dat is ideaal voor een transistor, maar ongeschikt voor een thermometer. Commerciële NTC's gebruiken meestal metaaloxiden met een geschikte band gap.
\end{oplossing}
